\section{Threat Model and Assumptions}
\label{sec:threat}
There are $N$ enrolled client accounts. Some control unique identities; others
may represent duplicate enrollment attempts. A trusted external registrar models
identity uniqueness with configurable false-accept and false-reject rates.
The prototype does not implement biometric, social-graph, or other
Proof-of-Personhood verification. Independently enrolled colluding participants
therefore remain in scope even when enrollment makes no errors.

The adversary can poison local updates and independently compromise validators.
Corrupted validators may approve poisoned updates, reject honest updates, or
report inflated monitoring quality. The experiment driver knows attacker labels
for injection and scoring. Honest validation, eligibility, aggregation, and
recovery do not receive these labels. ALIE experiments assume access to honest
reference updates and are labeled as an omniscient-adversary setting.

The registrar, off-chain experiment orchestrator, model aggregation computation,
and checkpoint store are trusted. The ledger authenticates account transactions
and enforces its contract state machine. It does not prove that a reported score
was obtained from an authentic validation partition. In particular, a malicious
majority or a biased committee can defeat the quality monitor.

Client data is divided into disjoint training and private validation partitions.
Held-out test examples are used exclusively for reporting. Trusted server
reference examples are reserved separately for methods that require them, and
are removed from all client partitions. Raw validation examples remain local;
model updates and released metrics may still leak information. Secure aggregation,
differential privacy, and a production identity service are outside this
implementation's guarantees.
